% =====================================================================
\chapter{Design: forward compression}\label{ch:design}
% =====================================================================

If comprehension is inverse compression, design is forward compression: compression
boundaries are built in advance for future readers. This chapter turns that symmetry into
an optimisable objective function, states its dynamical constraint, and shows where
layering fails.

\section{Read-sets and the objective function}\label{sec:readset}

The measurable objective of architecture is not an adjective such as ``decoupling'' or
``high cohesion'' but a number that can be counted.

\begin{definition}[Change read-set]\label{def:readset}
Let the code base consist of a set of modules $\Mod$, namely
$\bigcup_{c \in \mathcal{C}} \rho(c)$, where $\mathcal{C}$ is the set of changes. For a
change $c$, its \term{read-set} $\rho(c) \subseteq \Mod$ is the set of modules that must be
read in order to complete $c$ correctly. The \term{architectural quality metric} of the
system is
\begin{equation}\label{eq:readset}
Q(\Mod) \;=\; \mathbb{E}_{c \sim \mathcal{C}}\bigl[\,|\rho(c)|\,\bigr],
\end{equation}
that is, the mean read-set size weighted by the historical change distribution.
\end{definition}

Definition \ref{def:readset} is operational: it can be estimated directly from version
history, by counting for each merged change the number of top-level modules, files and
directories it touches and taking the resulting distribution. This is not a perfect proxy
(touching is not having to read), but it is \emph{measurable}, and ``degree of
modularity'' is not.

\begin{definition}[Enforcement sites of an invariant]\label{def:enforce}
Let $\Inv$ be the set of invariants of the system. For $\iota \in \Inv$, its
\term{enforcement-site set} $\mathrm{S}(\iota) \subseteq \Mod$ is the set of all modules
that \emph{must be implemented correctly for $\iota$ to be guaranteed}. Enforcement is
\term{effective} if $\iota$ is not violated in any execution. Its \term{unattested part}
$\mathrm{S}_{\mathrm{un}}(\iota) \subseteq \mathrm{S}(\iota)$ is the subset of those sites whose
correctness cannot be established without reading them: sites with no machine-checked proof, no
enforced type-level constraint, no separately checkable contract and no accepted certificate. Sites
outside $\mathrm{S}_{\mathrm{un}}(\iota)$ are still \emph{necessary}---if they were wrong the
invariant would fail---but they are \emph{substitutable by their evidence}, in the sense that an
agent may rely on the guarantee without inspecting the implementation.
\end{definition}

\begin{theorem}[Centralisation of invariants]\label{th:centralise}
Suppose that enforcement of invariant $\iota$ is effective and that a change $c$ may break
$\iota$, that is, there exists a legal sequence of changes that breaks $\iota$ in the
absence of an understanding of $\mathrm{S}(\iota)$. Then
\begin{equation}\label{eq:central}
|\rho(c)| \;\ge\; \bigl|\mathrm{S}_{\mathrm{un}}(\iota)\bigr| .
\end{equation}
Hence, other things equal, minimising $\mathbb{E}[|\rho|]$ requires minimising
$\sum_{\iota \in \Inv} |\mathrm{S}_{\mathrm{un}}(\iota)|$, whose lower bound is attained when every
enforcement site is either the only one or independently checkable.
\end{theorem}

\begin{proof}
If $c$ may break $\iota$, an agent executing $c$ must be able to judge whether the change is safe.
For each site in $\mathrm{S}_{\mathrm{un}}(\iota)$ that judgement requires reading that site, since
by definition no other evidence settles it; hence
$\rho(c) \supseteq \mathrm{S}_{\mathrm{un}}(\iota)$, which gives \eqref{eq:central}. Sites in
$\mathrm{S}(\iota) \setminus \mathrm{S}_{\mathrm{un}}(\iota)$ may be relied upon through their
evidence and need not be read. Summing over $\iota$:
$\mathbb{E}[|\rho|] \ge \mathbb{E}[|\bigcup_\iota \mathrm{S}_{\mathrm{un}}(\iota)|]$, with equality
in the optimistic case in which those sets are disjoint. The bound is attained at
$|\mathrm{S}_{\mathrm{un}}(\iota)| = 0$, that is, when every enforcement site carries evidence that
can be checked without being read---a type system, a private constructor, a capability token, a
database constraint or a proof.
\end{proof}

The change from $|\mathrm{S}(\iota)|$ to $|\mathrm{S}_{\mathrm{un}}(\iota)|$ is not a weakening
introduced for safety. It is the correction of a genuine conflation, and the distinction it draws
was missing from the argument. Three different sets were being treated as one: the modules whose
correct behaviour is necessary for $\iota$ (that is $\mathrm{S}(\iota)$); the evidence needed to
show a given change safe; and the source files an agent must actually read. The first does not bound
the third except through the second, and the second is exactly what $\mathrm{S}_{\mathrm{un}}$
isolates. A module that has a machine-checked proof or a stable, independently checkable interface
does not enter the read-set, no matter how necessary its correctness is.

Two consequences follow, and they are the practically useful ones. First, the bound explains why
encoding invariants in types has engineering value beyond expressiveness, and does so quantitatively:
it lowers $|\mathrm{S}_{\mathrm{un}}(\iota)|$ rather than merely relocating the invariant. Second, it
tells a maintainer where the read-set actually comes from: not from the number of places an invariant
matters, but from the number of those places whose correctness rests on nothing checkable. That is a
smaller and more actionable number, and it is the one a refactoring should target.

\begin{corollary}[Optimal location of enforcement sites]\label{cor:site}
Among all candidates satisfying the effectiveness constraint, the enforcement sites that minimise
the expected read-set of \eqref{eq:central} are the modules of \emph{least fan-in}, that is, those on
which fewest other modules depend and which are therefore least likely to be circumvented by
modification. The scope of this optimality is the read-set objective alone; see the remark below.
\end{corollary}

\begin{remark}[Limits of Theorem \ref{th:centralise}]\label{rmk:central-limit}
The theorem has one premise: enforcement must be \emph{effective}. In practice,
concentrating an invariant in a single place often fails, because callers can circumvent
that place. This is the real reason why ``scattered invariants'' recur in engineering. It
is not ignorance but a defence against the risk of circumvention. The correct reading of
Theorem \ref{th:centralise} is therefore: \emph{scattering invariants is a defence with a
cost; where invariants can genuinely be centralised (for instance by a type system, private
constructors, capability tokens or database constraints that make circumvention impossible),
centralisation is strictly better than scattering.}
This explains why ``encoding invariants in types'' has engineering value beyond
expressiveness: it lowers $|\mathrm{S}_{\mathrm{un}}(\iota)|$ towards $0$ without increasing the risk
of circumvention. The two mechanisms are independent and both matter---types remove the need to read,
effective enforcement removes the need to trust.
\end{remark}

\begin{remark}[Centralisation optimises one objective among several]\label{rmk:central-multi}
Corollary \ref{cor:site} identifies the read-set-optimal enforcement site. It should not be read as
identifying the best site, because the read-set is one objective among several, and the others are
not represented in the model at all. Moving an invariant into a single module also moves \emph{load}
onto that module, adds a hop to every operation that would otherwise be local, concentrates
\emph{failure impact}---if the enforcement site fails, every dependent fails together---and creates a
\emph{change bottleneck}, since every change touching the invariant now serialises through one
owner.

These are the same dimensions the paper insists must not be collapsed elsewhere
(Principle \ref{pr:crosscutting}), and the honest statement of the corollary is therefore
conditional: \emph{given} that effectiveness is achievable at the site in question and \emph{given}
that the other four dimensions are acceptable there, least fan-in is optimal for the read-set. When
they are not acceptable, the correct problem is a multi-objective one---correctness, change impact,
latency, and failure domain---and the read-set bound is a constraint within it rather than the
objective. A design review that treats ``centralise the invariant'' as an unconditional
recommendation is applying the corollary outside its scope.
\end{remark}

\section{Near-decomposability}\label{sec:neardecomp}

Simon's near-decomposability \cite{simon1962} is the classical explanation of why a level
structure is stable: interaction within a layer is dense, interaction between layers is
sparse. But ``sparse interaction'' is not itself testable. This section gives a testable
equivalent characterisation.

\begin{definition}[Near-decomposable]\label{def:neardecomp}
Let module $M$ have an interface $I$. For an execution $e$ of $M$, let
$\mathrm{trace}_I(e)$ be the sequence of interactions of $e$ on $I$ and
$\mathrm{behav}_{\mathrm{out}}(e)$ the behaviour of $e$ observable outside $M$. Then $M$ is
\term{near-decomposable} with respect to $I$ if there exists a function $f$ such that
\begin{equation}\label{eq:neardecomp}
\forall e:\quad \mathrm{behav}_{\mathrm{out}}(e) \;=\;
f\bigl(\mathrm{trace}_I(e)\bigr),
\end{equation}
that is, the observation on $I$ is a \emph{sufficient statistic} for the external behaviour
of $M$.
\end{definition}

The relation between Definition \ref{def:neardecomp} and Definition \ref{def:sufficient}
($T$-sufficiency) is immediate: take the query family $Q$ of a task $T$ to be
``observations of external behaviour''; near-decomposability is then exactly the
$T$-sufficiency of $M$ with respect to its interface. Simon's qualitative principle is
thereby translated into a decidable condition.

\begin{proposition}[Leak pairs]\label{prop:leakpair}
$M$ is \emph{not} near-decomposable with respect to $I$ if and only if there exist two
executions $e_1, e_2$ such that
\begin{equation}\label{eq:leakpair}
\mathrm{trace}_I(e_1) = \mathrm{trace}_I(e_2)
\qquad\text{and}\qquad
\mathrm{behav}_{\mathrm{out}}(e_1) \ne \mathrm{behav}_{\mathrm{out}}(e_2).
\end{equation}
Such a pair is called a \term{leak pair}.
\end{proposition}

\begin{proof}
$(\Leftarrow)$ If such $e_1, e_2$ exist, then no function $f$ can satisfy both
$f(\mathrm{trace}_I(e_1)) = \mathrm{behav}_{\mathrm{out}}(e_1)$ and
$f(\mathrm{trace}_I(e_2)) = \mathrm{behav}_{\mathrm{out}}(e_2)$, since the arguments on the
two sides are equal while the assigned values are not; hence \eqref{eq:neardecomp} fails.
$(\Rightarrow)$ If \eqref{eq:neardecomp} fails, no $f$ satisfies the condition, and by the
definition of a function there must exist two points with the same argument and different
assigned values, which is \eqref{eq:leakpair}.
\end{proof}

\begin{corollary}[A diagnostic procedure]\label{cor:diagnose}
Proposition \ref{prop:leakpair} yields an \emph{executable} diagnostic: construct two
scenarios identical at the interface but differing in internal state or external
environment, and observe whether the external behaviour diverges. If it diverges, the
interface is not a sufficient statistic, and the module boundary leaks there.
\end{corollary}

Corollary \ref{cor:diagnose} turns ``does this module's abstraction leak'' from an aesthetic
judgement into a differential test. Its practical value is that \emph{whether an abstraction
leaks is measurable, not merely felt}. The engineering commonplaces of ``shadow state'',
``implicit globals'' and ``cross-request cache pollution'' are all concrete instances of
leak pairs.

\begin{principle}[An equivalent design objective]\label{pr:design-eq}
By Corollary \ref{cor:diagnose}, the testable objective of design can be written as:
\emph{for every claimed module boundary, the interface observation should be a sufficient
statistic for external behaviour; the existence of a leak pair is the failure of that
boundary.}
\end{principle}

\begin{principle}[Two criteria for a claimed boundary, not one]\label{pr:two-criteria}
A claimed boundary between a module and its environment makes two separate claims, and they fail
separately.
\begin{enumerate}[label=(\alph*)]
\item \term{Sufficiency}: the interface observation determines the external behaviour. This is
      Definition \ref{def:neardecomp}; its failure is a leak pair (Proposition \ref{prop:leakpair}),
      found by differential testing (Corollary \ref{cor:diagnose}).
\item \term{Completeness}: the observation map the boundary declares is the whole of what is
      observable across it. Its failure is a difference the boundary declared irrelevant becoming
      visible anyway---through timing, resource occupancy or failure behaviour---and it is found by
      the techniques of information-flow analysis rather than by differential testing
      (Remark \ref{rmk:noninterference}).
\end{enumerate}
The two are independent, in both directions. A boundary can be sufficient and incomplete: every pair
of interface-identical executions has the same external behaviour, and yet a declared-irrelevant
difference is readable from above through a timing channel. It can equally be complete and
insufficient: nothing leaks through undeclared channels, while the declared interface is too coarse
to determine the behaviour.

Corollary \ref{cor:diagnose} and Principle \ref{pr:design-eq} address the first criterion only.
The second is the one that fails silently, because a boundary that leaks through an undeclared
channel passes every differential test written against the declared interface---the test never
constructs the offending pair, since the pair differs in something the boundary said did not matter.
\end{principle}

\section{Technical debt}\label{sec:debt}

``Technical debt'' is usually a metaphor in engineering usage. This section gives a
definition that is precise and still useful.

\begin{definition}[Technical debt]\label{def:debt}
Given an encoding, let the description length of the system code be $\DL(\text{code})$ and
let its \term{shortest generating description} be
$\DL^{\star} = \min\{\,|d| : d \text{ generates code}\,\}$, that is, the shortest program
that generates that code in the Kolmogorov sense. Define the \term{technical debt}
\begin{equation}\label{eq:debt}
D \;=\; \DL(\text{code}) \;-\; \DL^{\star}.
\end{equation}
\end{definition}

The definition is defensible but non-computable ($\DL^{\star}$ is not computable), so its
value lies in the \emph{direction} and the \emph{criterion} it supplies, not in a numerical
value.

\begin{principle}[How technical debt accrues and is repaid]\label{pr:debt}
Technical debt grows when \emph{a change is not re-abstracted}: if a modification is not
re-abstracted, the gap between the code and its ``shortest generating description'' widens.
Therefore:
\begin{enumerate}[label=(\alph*)]
\item refactoring $=$ recompression, a survival necessity for holding the compression ratio
      above the comprehensibility threshold;
\item a system that is not refactored does not stop growing; it accumulates
      incomprehensibility;
\item deciding whether a refactoring is worthwhile is equivalent to deciding whether it
      significantly lowers $D$, that is, whether it makes the system easier to generate
      from a shorter description.
\end{enumerate}
\end{principle}

Principle \ref{pr:debt}(c) gives a practical criterion: if a refactoring does not let you
say in fewer words what the module does, it is probably only reformatting.

\begin{remark}[A rate--distortion reading of refactoring]\label{rmk:rd-refactor}
Regard $(\DL(M), L_T(M))$ as a point in the rate--distortion plane \cite{shannon1959}
(Definition \ref{def:mdl}). The current state of a system is a point on the curve or outside
it: \emph{outside the curve} means that a model of the same description length attains lower
loss, that is, that there is pure room for improvement; refactoring is then moving along the
curve or moving the point back onto it.

This reading yields a testable prediction: \emph{the slope of the rate--distortion curve,
$\partial L_T / \partial \DL$, should predict which refactorings pay off.} A refactoring at
a point where the slope is zero does not lower task loss and is therefore not worth doing;
where the slope is large, it is worth doing even at high cost. If that predictive power does
not hold in real projects, the choice of objective function in this chapter requires
revision (Chapter \ref{ch:program}, hypothesis $H_{12}$).
\end{remark}

\section{Governable complexity and comprehension capacity}\label{sec:governable}

Section \ref{sec:readset} defines architectural quality as a static quantity
$\mathbb{E}[\rho]$. Real systems, however, have a \emph{rate}: changes arrive at some
frequency, and an organisation's reading capacity is finite. This section supplies that
dynamics; it is the only place in the paper where time is a variable.

\begin{definition}[Governable complexity and comprehension capacity]\label{def:governable}
Let $\lambda(t)$ be the change arrival rate (changes per unit time) and let
$\mathbb{E}[\rho]$ be the expected read-set (Definition \ref{def:readset}). Define the
\term{governable complexity}
\begin{equation}\label{eq:governable}
C(t) \;=\; \mathbb{E}\bigl[\rho\bigr] \cdot \lambda(t),
\end{equation}
that is, the information throughput required to maintain the system, in units of
module-reads per unit time. Let the \term{comprehension capacity} of the organisation be
\begin{equation}\label{eq:capacity}
K \;=\; h \cdot b,
\end{equation}
where $h$ is the number of people who can be applied to comprehension and $b$ is the
comprehension bandwidth of one person (module-reads per unit time). The system is in a
\term{governable state} if $C(t) \le K$ for all relevant $t$.
\end{definition}

\begin{proposition}[The scaling exponent of required headcount, under the capacity model]\label{prop:exponent}
Let a system have $n$ features, each changing independently at rate $\lambda_0$, so that
$\lambda \propto n$. Let the read-set $\mathbb{E}[\rho]$ scale as $n^{\alpha}$ in $n$, with
$\alpha \ge 0$ a parameter of the model, $\alpha = 0$ corresponding to complete
localisation and $\alpha = 1$ to unconstrained coupling. Then the governable complexity
scales as
\begin{equation}\label{eq:exponent}
C \;\propto\; n^{1+\alpha}.
\end{equation}
From $C \le K$ and $K = hb$ with $b$ constant, the headcount needed to keep the system governable
scales as $h \propto n^{1+\alpha}$. Therefore, \emph{within this model}:
\begin{enumerate}[label=(\alph*)]
\item if $\alpha$ is pressed to $0$, headcount grows linearly with system size and the system can
      persist indefinitely;
\item if $\alpha = 1$, headcount must grow as the \emph{square} of size, and the system must
      become unstable at some size.
\end{enumerate}
\end{proposition}

\begin{proof}
By Definition \ref{def:governable}, $C = \mathbb{E}[\rho]\cdot\lambda$. Substituting
$\lambda = n\lambda_0$ and $\mathbb{E}[\rho] = \Theta(n^{\alpha})$ gives
$C = \Theta(n^{1+\alpha})$. From $C \le K = hb$ it follows that
$h \ge C/b = \Theta(n^{1+\alpha})$.
\end{proof}

Two things the proposition does not establish, and that the earlier title---\emph{architecture fixes
the scaling exponent}---asserted by implication. First, it does not show that the exponent is fixed
by the architecture; it \emph{assumes} that $\mathbb{E}[\rho]$ scales as a power of $n$ and treats
$\alpha$ as the parameter of that scaling. Whether $\alpha$ is architecture-determined, and whether
the read-set proxy used to estimate it is valid, are both empirical questions---the second is
Hypothesis $H_4$, and the first is what Hypothesis $H_{14}$ asks directly. Second, it does not show
that real organisations' headcount follows the exponent. It shows that headcount must follow it
\emph{if} the organisation is to remain governable under the stated capacity model, which is a
conditional statement about a model of comprehension capacity, not a forecast about hiring. The
empirical content is that systems with higher estimated $\alpha$ should require disproportionately
more coordination effort per unit of growth; that is testable, and Chapter \ref{ch:program} proposes
how.

\begin{corollary}[The actual role of architecture]\label{cor:arch-role}
Proposition \ref{prop:exponent} assigns architecture a precise duty: it does not lower the
cost of an individual change but \emph{lowers the exponent with which cost grows with size}.
This explains why architectural benefit is invisible in small systems and decisive for
survival in large ones, and why the charge of ``architectural over-engineering'' always
appears to hold at small scale.
\end{corollary}

\begin{corollary}[Derivation of two classical conclusions]\label{cor:brooks}
Equation \eqref{eq:exponent} yields two classical conclusions:
\begin{enumerate}[label=(\alph*)]
\item \term{Relation between organisation size and system size}: from
      $h \propto n^{1+\alpha}$, organisation size must grow with system size, and the
      growth exponent is fixed by the architecture. This is the \emph{dynamical} version of
      the correspondence between organisational structure and system structure
      \cite{conway1968}; that correspondence is a static statement (structure mirrors
      organisation), whereas the dynamical version gives the governability bound.
\item \term{Adding people makes a project slower}: increasing $h$ raises $K$ linearly, but
      at the same time raises the number of changes in flight and hence $\lambda$, and
      raises $\mathbb{E}[\rho]$ through the interfaces it adds. When the growth of the
      latter two exceeds that of the former, $C/K$ worsens rather than improves. This is an
      in-model explanation of the empirical rule that adding people to a late project makes
      it later.
\end{enumerate}
\end{corollary}

\begin{principle}[Design implication]\label{pr:governable-design}
By Proposition \ref{prop:exponent}, the criterion of an architecture review should not be
``is the read-set small'' but \emph{``what is the scaling exponent of the read-set in
size''}. A system whose read-set is currently small but grows rapidly with size is more
dangerous than one whose read-set is currently larger but scales flatly. This requires
hypothesis $H_4$ to be strengthened: the value of a read-set lies mainly in its
\emph{scaling}, not in its \emph{magnitude}.
\end{principle}

\begin{remark}[The only dynamics]\label{rmk:only-dynamics}
The paper has so far been static: structure, invariants, read-sets and contracts are all
time-invariant quantities. Definition \ref{def:governable} is the only dynamical variable,
because it is the only place where ``what an organisation can process'' and ``what the
system demands'' must match in \emph{rate}. Elsewhere time can be eliminated.
\end{remark}

\section{The design procedure}\label{sec:seven-steps}

\begin{enumerate}[label=(D\arabic*)]
\item \term{Fix the invariants first}. Before writing any code, write down the propositions
      that must hold forever of this system (for example ``the total balance is
      conserved'', ``an I/O thread never blocks'', ``the version number increases
      monotonically''). Invariants are the gravitation of a level; all structure
      crystallises around them. By Theorem \ref{th:centralise}, each invariant should
      designate a single enforcement site and state why circumvention is impossible.
\item \term{Choose the primitives}. Atoms, components, messages, transactions, entities,
      and so on. Once the primitives are chosen they become the shared language of all
      subsequent discussion of the system. \emph{Choosing the primitives is the most
      important single decision at this level}, because it fixes the shape of every
      expression thereafter.
\item \term{Fix the data model}. Entities, relations, lifecycles, ownership. A data model
      is an order of magnitude more stable than logic, so it should be frozen first.
\item \term{Cut the organs}. The basis of modularisation is not ``similar
      functionality'' but \emph{information that changes together being hidden together}
      \cite{parnas1972}. The test for each module is that its interface is small enough to
      hold in mind at once and its implementation deep enough to be replaced wholesale
      \cite{ousterhout2018}. The interface is the compression summary you deliberately
      maintain; the implementation is the compressed content that may be rewritten.
\item \term{Make the physical laws executable}. An invariant written in a document rots;
      only encoded in a type system, assertions, contract tests and property tests does it
      endure. \emph{Types are architectural documentation enforced by the compiler}. This
      is the confluence of ``designing for comprehensibility'' and ``designing for
      correctness'': an invariant that a machine can check need not be maintained from
      memory by a human or an AI.
\item \term{Record decisions}. Write down why each choice was made and what was given up.
      Code is the result; a decision record is part of the generating rules, and it spares
      the reader performing inverse compression the most expensive guessing.
\item \term{Refactoring $=$ recompression}. See Principle \ref{pr:debt}.
\end{enumerate}

One meta-rule runs through all seven:

\begin{principle}[Design for the reader's mental bandwidth]\label{pr:reader}
Design for the mental bandwidth of your reader, not for your typing speed. Code is read far
more often than it is written, and the users of an architecture are future comprehenders.
\end{principle}

\section{Cross-cutting concerns and the limits of layering}\label{sec:crosscutting}

The failure modes of layering must be acknowledged explicitly. One class of concerns cannot
be localised to any layer:

\begin{itemize}
\item security and authorisation (every module may introduce a vulnerability);
\item multi-tenant isolation (spanning storage, computation and network);
\item distributed tracing (which must run through every call path);
\item transactions and consistency (across the boundaries of state owners);
\item observability and quotas (required at every point of resource consumption).
\end{itemize}

\begin{principle}[Cross-cutting concerns are independent dimensions]\label{pr:crosscutting}
For cross-cutting concerns, the mature treatment is to acknowledge them as
\emph{independent dimensions} (policy plane, capability model, taint propagation, service
mesh) rather than forcing them into the hierarchy and pretending that they are a layer.
\end{principle}

This principle follows directly from the discussion in Chapter \ref{ch:convergence}: a
level structure is supported by near-decomposability, and near-decomposability presupposes
sparse interaction between layers. When a concern has a non-trivial interaction with every
module, it breaks near-decomposability by definition, and therefore cannot be represented as
``a layer''. The result of forcing the representation is that the level structure fails, and
the failure is usually misdiagnosed as ``insufficiently thorough layering''.

\section{Design checklist}\label{sec:checklist}

The following checklist turns this chapter into executable items. Each item should be
answered explicitly at a design review.

{\small\begin{longtable}{L{1.0cm}L{14.0cm}}
\caption{Design checklist for large systems}\label{tab:checklist}\\
\toprule
\textbf{No.} & \textbf{Check} \\
\midrule
\endfirsthead
\toprule
\textbf{No.} & \textbf{Check} \\
\midrule
\endhead
\bottomrule
\endlastfoot
A1 & Which propositions must hold forever of the system? Where is each enforced? Why is circumvention impossible?\\
A2 & Is the enforcement site of each invariant unique? If not, has the cost (the read-set) been explicitly accepted?\\
A3 & What are the primitives of this level? Are they the shared language of all subsequent discussion?\\
A4 & Has the data model been frozen before the logic? Are entities, relations, lifecycles and ownership explicit?\\
\midrule
B1 & Is each module's interface small enough to hold in mind at once?\\
B2 & Is each module deep enough to be replaced wholesale?\\
B3 & For each claimed boundary, does a leak pair exist? (That is: two executions with the same interface trace and different external behaviour.)\\
B4 & Has each module's \emph{oblivious set} been written down explicitly?\\
B5 & Which modules change together? Has this been verified against co-change history rather than intuition alone?\\
\midrule
C1 & How many unrelated modules must be modified to add a new feature? (The measured read-set.)\\
C2 & Does replacing the underlying implementation force large-scale modification of the business logic?\\
C3 & When a subsystem fails, where does the failure propagate?\\
C4 & Does adding a task type require rewriting the shared scheduling logic?\\
C5 & Does concurrent access break state invariants? Where is this detected?\\
C6 & Are performance bottlenecks concentrated in a few critical paths?\\
\midrule
D1 & Have the invariants been encoded in executable form (types, assertions, contract tests, property tests)?\\
D2 & Can each module's contract be verified in isolation? If not, is it a module?\\
D3 & Does each critical subsystem have a traceable record of what is known about it (Table \ref{tab:record})?\\
D4 & Does the decision record cover why each choice was made and what was given up?\\
D5 & Are cross-cutting concerns treated as independent dimensions rather than disguised as a layer?\\
\midrule
E1 & Does the system need a unified substrate? Judge by $n \times m$, not by lines of code.\\
E2 & Of the six functions, which are built and which are adopted? The criterion is the domain-criticality of the invariants.\\
E3 & For the parts that are built, are their invariants ones that a general-purpose host cannot provide?\\
E4 & For the parts that are adopted, is the strength of the guarantee offered by the host sufficient to support this layer's contract?\\
\bottomrule
\end{longtable}}

Theorem \ref{th:centralise} and Propositions \ref{prop:leakpair} and \ref{prop:exponent}
hold under the stated models. The empirical content is hypotheses $H_4$ and $H_{12}$: that
the scaling exponent of the read-set outperforms read-set magnitude and conventional size
metrics; that centralising invariants (where circumvention is impossible) lowers both the
read-set and the defect rate; and that the rate--distortion slope predicts the benefit of a
refactoring. If all three are refuted, the choice of objective function in this chapter
requires revision.
